\section{An exact geometric collision formula and harmonic hierarchy laws}
\label{hc:section}

The incoming mixed-spectral report already determines the collision of
three undifferentiated Stieltjes factors at shifts $(0,a\epsilon,b\epsilon)$
for fixed $0<a<b$. It also supplies the complete spectral subtraction
at a fixed, arbitrarily multiple collision \cite{RepoMixed}. These are
substantial antecedents. A moving configuration such as
$(0,\epsilon,\epsilon+\epsilon^2)$ is a different question: inserting
$\epsilon$-dependent rates in a fixed-rate error estimate is not justified
without uniform control. This section proves an exact formula before
any rate specialization. It permits arbitrarily many polygamma factors,
arbitrary argument-derivative orders, and analytic collision arcs with
several different vanishing orders. Its correction is a finite expression
in polygamma values, derivatives, rational functions, and logarithms.
The corrected remainder is holomorphic in all the shifts.

\subsection{The coordinate convention and the local telescoping identity}

Write
\[
 f_r(z)=-\psi^{(r)}(z),\qquad
 n_i=r_i+1,\qquad \kappa_i=(-1)^{r_i}r_i!,\qquad
 h_i(z)=f_{r_i}(1+z).
\]
Thus $f_0=\gamma_0$ and $f_r=\gamma_0^{(r)}$; this section concerns
Stieltjes index zero with arbitrary argument derivatives. In particular,
no claim about all positive Stieltjes indices is hidden in the notation.
The polygamma recurrence \cite{DLMFPolygamma} gives the exact local formula
\begin{equation}\label{hc:one-sided}
 f_{r_i}(\{x\})=h_i(x)+
       \kappa_i\mathbf1_{x>0}x^{-n_i},\qquad 0<|x|<\delta<1.
\end{equation}
The functions $h_i$ are holomorphic for $|x|<1$.

For real shifts $a_1<\cdots<a_d$ sufficiently close to zero, define
\begin{equation}\label{hc:Cdef}
 \mathcal C_{\boldsymbol r}(\boldsymbol a)
 =\FP_x\int_0^1\prod_{i=1}^d f_{r_i}(\{x+a_i\})\,\dd x,
 \qquad
 Q_{\boldsymbol r}=\FP_x\int_0^1\prod_{i=1}^d f_{r_i}(x)\,\dd x.
\end{equation}
At the seam $x=-a_i$ on the circle, the cutoff coordinate is exactly
$x+a_i$, taken positive, and the deleted right interval has ambient
length $\eta_i$. The finite part is the coefficient independent of all
$\eta_i$ and $\log\eta_i$. The cutoffs may be taken independently.
These are pointwise products away from finitely many distinct seams;
there is no multiplication of distributions with overlapping singular
supports. The merged quantity uses the same ambient coordinate $x$.

Put, for $1\le m\le d$,
\begin{equation}\label{hc:Km}
 K_m(z,\boldsymbol a)=
 \frac{\kappa_m}{(z+a_m)^{n_m}}
 \prod_{k<m}h_k(z+a_k)
 \prod_{k>m}f_{r_k}(z+a_k).
\end{equation}
In a fixed neighborhood of the merged seam one has
\begin{equation}\label{hc:telescope}
 \prod_{i=1}^d f_{r_i}(\{x+a_i\})
 =\prod_{i=1}^d h_i(x+a_i)
      +\sum_{m=1}^d\mathbf1_{x>-a_m}K_m(x,\boldsymbol a).
\end{equation}
Indeed, expand the product according to the first factor for which the
singular part in \eqref{hc:one-sided} is selected. If this first index is
$m$ and $x>-a_m$, then $x>-a_k$ for every $k>m$, so all the subsequent
factors may be replaced by their meromorphic germs $f_{r_k}(x+a_k)$.
This telescoping form replaces a sum over $2^d$ subsets by $d$ terms.

For $m<i$ and $1\le p\le n_i$, let $B_{mi,p}(\boldsymbol a)$ be the
coefficient of $(z+a_i)^{-p}$ in the principal part of $K_m$ at $-a_i$.
It is explicitly
\begin{align}
 B_{mi,p}(\boldsymbol a)
 =\frac{\kappa_i}{(n_i-p)!}
 \left.\partial_z^{\,n_i-p}\left\{
 \frac{\kappa_m}{(z+a_m)^{n_m}}
 \prod_{k<m}h_k(z+a_k)
 \prod_{\substack{k>m\\k\ne i}}f_{r_k}(z+a_k)
 \right\}\right|_{z=-a_i}.
 \label{hc:B}
\end{align}
Only finitely many derivatives are required. All arguments after
this substitution are differences of shifts.

\subsection{Exact subtraction, with a jointly holomorphic remainder}

\begin{theorem}[Geometric collision formula]\label{hc:main}
With the convention of \eqref{hc:Cdef}, set
\begin{equation}\label{hc:counter}
 \mathcal A_{\boldsymbol r}(\boldsymbol a)=
 \sum_{m<i}\left\{
 -B_{mi,1}(\boldsymbol a)\log(a_i-a_m)
 +\sum_{p=2}^{n_i}
  \frac{B_{mi,p}(\boldsymbol a)}{(p-1)(a_i-a_m)^{p-1}}
 \right\}.
\end{equation}
There is a function $\mathcal R_{\boldsymbol r}$ holomorphic in all
$a_i$ on a complex neighborhood of $\boldsymbol a=0$ such that, in
the stated real ordering chamber,
\begin{equation}\label{hc:exact}
 \boxed{\quad
 \mathcal C_{\boldsymbol r}(\boldsymbol a)
 =\mathcal A_{\boldsymbol r}(\boldsymbol a)
       +\mathcal R_{\boldsymbol r}(\boldsymbol a),\qquad
 \mathcal R_{\boldsymbol r}(\boldsymbol0)=Q_{\boldsymbol r}.
 \quad}
\end{equation}
The theorem applies to every number of factors and every tuple of
nonnegative derivative orders. It supplies the complete collision
correction, including its finite terms, rather than just its leading
divergences.
\end{theorem}

\begin{proof}
Choose a fixed small $\delta>0$ and represent the circle near its seam
by $x\in(-\delta,\delta)$. The integral outside this neighborhood and
the integral of the all-$h_i$ term in \eqref{hc:telescope} are
holomorphic in the shifts. Consider one $K_m$ term, integrated from
$-a_m$ to $\delta$ with its right endpoint cutoff at $-a_m$.

It is important not to take the limits of its individual principal-part
coefficients, which can diverge when poles coalesce. Define
\[
 D_m(z,\boldsymbol a)=\prod_{i=m}^d(z+a_i)^{n_i},\qquad
 H_m(z,\boldsymbol a)=D_m(z,\boldsymbol a)K_m(z,\boldsymbol a).
\]
The numerator is jointly holomorphic: explicitly it is
\[
 H_m=\kappa_m\prod_{k<m}h_k(z+a_k)
       \prod_{k>m}\bigl\{\kappa_k+(z+a_k)^{n_k}h_k(z+a_k)\bigr\}.
\]
For distinct shifts, write the polar decomposition
\[
 K_m(z,\boldsymbol a)
   =\sum_{i\ge m}\sum_{p=1}^{n_i}
          \frac{B_{mi,p}(\boldsymbol a)}{(z+a_i)^p}
       +W_m(z,\boldsymbol a),
\]
where the notation $B_{mm,p}$ is used here only for the analogous
principal coefficients at the starting point.
The remainder extends jointly holomorphically as all the shifts merge.
For example, on a fixed circle $|w|=\rho>\delta$ inside the domain of
all $H_m$, with the shifts sufficiently small, the residue theorem gives
\begin{equation}\label{hc:Hermite}
 W_m(z,\boldsymbol a)=\frac1{2\pi\ii}
 \int_{|w|=\rho}\frac{H_m(w,\boldsymbol a)}
  {D_m(w,\boldsymbol a)(w-z)}\,\dd w,\qquad |z|<\rho.
\end{equation}
The denominator stays away from zero on this contour, proving the
claimed joint holomorphy. This is the analytic cancellation usually
encoded by confluent Hermite divided differences.

A primitive of the polar part is
\[
 U_m(z,\boldsymbol a)=
 \sum_{i\ge m} B_{mi,1}\log(z+a_i)
 -\sum_{i\ge m}\sum_{p=2}^{n_i}
       \frac{B_{mi,p}}{p-1}(z+a_i)^{1-p}.
\]
At the upper point $z=\delta$ its apparent collision poles cancel.
To see this without a coefficientwise limit, choose a second, smaller
circle $|w|=\rho_0<\delta$ enclosing all $-a_i$. Then
\begin{equation}\label{hc:upper}
 U_m(\delta,\boldsymbol a)=\frac1{2\pi\ii}
  \int_{|w|=\rho_0}K_m(w,\boldsymbol a)
                         \log(\delta-w)\,\dd w.
\end{equation}
The logarithm is holomorphic inside this smaller circle. The displayed
integral is jointly holomorphic in the shifts. For $p\ge2$, the derivative
of the logarithm in the residue calculation is
$-(p-2)!/(\delta+a_i)^{p-1}$, giving exactly the sign in $U_m$.

At the lower point $z=-a_m+\eta$, the terms belonging to $i=m$ contain
only $\log\eta$ and negative powers of $\eta$; their cutoff constant
is zero. For $i>m$, ordinary substitution at $\eta=0$ is legitimate.
Subtracting these lower primitive values gives precisely the $m$th
summand in \eqref{hc:counter}. Consequently the remainder of the local
integral is
\[
 U_m(\delta,\boldsymbol a)
       +\int_{-a_m}^{\delta}W_m(z,\boldsymbol a)\,\dd z,
\]
which is jointly holomorphic, by \eqref{hc:Hermite} and
\eqref{hc:upper}. At zero shifts it is exactly the ambient finite part
of $K_m(z,\boldsymbol0)$ on $(0,\delta)$: all merged lower polar
primitive terms again have zero cutoff constant. Finally sum over $m$
and add the holomorphic all-regular and exterior integrals. The
zero-shift form of \eqref{hc:telescope} identifies their total with
$Q_{\boldsymbol r}$. This proves \eqref{hc:exact}.
\end{proof}

\begin{corollary}[Analytic collision arcs]\label{hc:arcs}
Let $a_i(\epsilon)$ be real analytic, vanish at zero, and be pairwise
distinct for sufficiently small $\epsilon>0$, with a fixed ordering
there. Then the complete expansion of
$\mathcal C_{\boldsymbol r}(\boldsymbol a(\epsilon))$ has the form
\[
 U(\epsilon)+V(\epsilon)\log\epsilon,
\]
where $U$ and $V$ are convergent Laurent germs with finite principal
parts. Its coefficient independent of $\epsilon$ and $\log\epsilon$ is
\begin{equation}\label{hc:arc-constant}
 \CT_{\epsilon}\mathcal C_{\boldsymbol r}(\boldsymbol a(\epsilon))
 =Q_{\boldsymbol r}
       +\CT_{\epsilon}\mathcal A_{\boldsymbol r}(\boldsymbol a(\epsilon)).
\end{equation}
Every term of the second constant is explicitly determined by finitely
many jets of the arcs and of the polygamma functions at $1$.
\end{corollary}

\begin{proof}
Each positive gap is
$a_i-a_m=b_{im}\epsilon^{\nu_{im}}(1+O(\epsilon))$, with
$b_{im}>0$ and positive integer $\nu_{im}$. Its logarithm is
$\nu_{im}\log\epsilon+\log b_{im}$ plus a convergent Taylor series.
Formula \eqref{hc:B} is meromorphic after composition with these arcs:
all possible poles come from a finite number of powers of nonzero
analytic gaps. Thus \eqref{hc:counter} has the stated form and its
constant needs only finitely many coefficients. The holomorphic
remainder contributes exactly its value at zero. No higher powers
of $\log\epsilon$ arise in this polygamma case.
\end{proof}

\subsection{A short formula at Stieltjes index zero}

Let $f=f_0=-\psi$ and $h(z)=f(1+z)$. If every $r_i=0$, all poles in
\eqref{hc:Km} are simple and \eqref{hc:counter} collapses to
\begin{equation}\label{hc:short}
 \boxed{\quad
 \mathcal A_{\boldsymbol0}(\boldsymbol a)=
 \sum_{m<i}\frac{\log(a_i-a_m)}{a_i-a_m}
   \prod_{k<m}h(a_k-a_i)
   \prod_{\substack{k>m\\k\ne i}}f(a_k-a_i).
 \quad}
\end{equation}
The negative arguments here are values of the meromorphic digamma
germ, not fractional parts. They are defined because the small gaps
are nonzero. Formula \eqref{hc:short} has only $d(d-1)/2$ summands.

For three factors at shifts $(0,a,b)$ it reads
\begin{equation}\label{hc:triple-exact}
 \mathcal A_3(0,a,b)=
 \frac{\log a}{a}f(b-a)
 +\frac{\log b}{b}f(a-b)
 +\frac{\log(b-a)}{b-a}h(-b).
\end{equation}
This exact identity is the key to the hierarchical limits below. For
fixed $a=\alpha\epsilon$, $b=\beta\epsilon$, expanding
\begin{equation}\label{hc:hseries}
 h(z)=\gamma+\sum_{k\ge1}(-1)^k\zeta(k+1)z^k
\end{equation}
recovers every displayed term through order $\epsilon^0$ in the
incoming fixed-rate triple theorem \cite{RepoMixed}. In particular,
for $(0,\epsilon,2\epsilon)$ the finite discrepancy is
$\zeta(2)\log2/2$, as already established there. The addition in
\eqref{hc:triple-exact} is the exact correction at arbitrary small
$a,b$, with a jointly holomorphic remainder, so rates may now coalesce.

Even a binary collision illustrates why the parameter matters. For
$d=2$, \eqref{hc:short} is $\log(a_2-a_1)/(a_2-a_1)$. Thus
\begin{equation}\label{hc:binary-curved}
 \CT_{\epsilon}\mathcal C_{0,0}(0,\epsilon+b\epsilon^2)
     =Q_{0,0}+b,
\end{equation}
because
$\log(\epsilon+b\epsilon^2)/(\epsilon+b\epsilon^2)
 =\epsilon^{-1}\log\epsilon-b\log\epsilon+b+O(\epsilon|\log\epsilon|)$.
This does not change the preceding binary theorem in its stated shift
coordinate. It records the effect of choosing a nonlinear collision arc.

\subsection{Two different hierarchies and their harmonic constants}

Set $Q_3=Q_{0,0,0}$, and write $H_n=\sum_{j=1}^n1/j$ for $n\ge1$.
In the formulas below a term containing $H_{u/v}$ is included only
when $v$ divides $u$; no value at a nonintegral index is needed.

\begin{theorem}[Oriented harmonic hierarchy laws]\label{hc:hierarchy}
Let $p>q\ge1$ be integers and $v=p-q$. The left nested collision obeys
\begin{equation}\label{hc:left}
 \boxed{\quad
 \CT_{\epsilon}\mathcal C_{0,0,0}(0,\epsilon^p,\epsilon^q)
   =Q_3-\mathbf1_{v\mid q}\,\gamma H_{q/v}.
 \quad}
\end{equation}
The right nested collision has the different value
\begin{equation}
 \boxed{\begin{aligned}
 \CT_{\epsilon}\mathcal C_{0,0,0}
           (0,\epsilon^q,\epsilon^q+\epsilon^p)
 =Q_3
 &+\mathbf1_{v\mid p+q}(-1)^{(p+q)/v}H_{(p+q)/v}
 \\[-2pt]
 &+\mathbf1_{v\mid q}\,\gamma(-1)^{q/v+1}H_{q/v}.
 \end{aligned}}\label{hc:right}
\end{equation}
In particular,
\begin{align}
 \CT_{\epsilon}\mathcal C_{0,0,0}(0,\epsilon^2,\epsilon)
     &=Q_3-\gamma,\label{hc:left-example}\\
 \CT_{\epsilon}\mathcal C_{0,0,0}(0,\epsilon,\epsilon+\epsilon^2)
     &=Q_3+\gamma-\frac{11}{6},\label{hc:right-example}\\
 \CT_{\epsilon}\mathcal C_{0,0,0}(0,\epsilon,\epsilon+\epsilon^3)
     &=Q_3+\frac32.\label{hc:right-second}
\end{align}
\end{theorem}

\begin{proof}
The two elementary, convergent generating functions needed are
\begin{equation}\label{hc:harmonic-generators}
 \frac{\log(1-z)}{1-z}=-\sum_{n\ge1}H_nz^n,
 \qquad
 \frac{\log(1+z)}{1+z}=\sum_{n\ge1}(-1)^{n+1}H_nz^n.
\end{equation}
The first follows by multiplying
$-\log(1-z)=\sum_{j\ge1}z^j/j$ by the geometric series; the second
is obtained by replacing $z$ with $-z$.

For the left nested configuration, substitute
$a=\epsilon^p$, $b=\epsilon^q$ into \eqref{hc:triple-exact}.
The first two logarithms are respectively $p\log\epsilon$ and
$q\log\epsilon$. They therefore make no contribution to the
log-free constant, regardless of the poles of their other factors.
The only possible contribution is
\[
 \epsilon^{-q}\frac{\log(1-\epsilon^v)}{1-\epsilon^v}
                       h(-\epsilon^q)
 =-\sum_{n\ge1}H_n\epsilon^{nv-q}
       \left\{\gamma+\sum_{k\ge1}\zeta(k+1)\epsilon^{qk}\right\}.
\]
For $k\ge1$, the exponent $nv-q+qk$ is strictly positive. The constant
comes solely from $k=0$ and $nv=q$, proving \eqref{hc:left} after
\eqref{hc:arc-constant} is used.

For the right nested configuration, put
$a=\epsilon^q$, $b=\epsilon^q+\epsilon^p$. Now the first and third
logarithms in \eqref{hc:triple-exact} are pure multiples of
$\log\epsilon$. The log-free part of the second term is
\begin{align*}
 &\frac{\log(1+\epsilon^v)}{1+\epsilon^v}
          \epsilon^{-q}f(-\epsilon^p)\\
 &\qquad=\left\{\sum_{n\ge1}(-1)^{n+1}H_n\epsilon^{nv}\right\}
       \left\{-\epsilon^{-p-q}+\gamma\epsilon^{-q}
                +\sum_{k\ge1}\zeta(k+1)\epsilon^{pk-q}\right\}.
\end{align*}
All terms in the last sum have positive exponents. The pole term
contributes only if $nv=p+q$, with coefficient $(-1)^nH_n$; the
constant regular term contributes only if $nv=q$, with coefficient
$\gamma(-1)^{n+1}H_n$. This proves \eqref{hc:right}. The examples use
$H_1=1$, $H_2=3/2$, and $H_3=11/6$.
\end{proof}

These two arrangements are not interchangeable. The original periodic
Stieltjes germ has a one-sided singularity, and reflection is not a
symmetry of its regular part. For example, subtracting
\eqref{hc:left-example} from \eqref{hc:right-example} eliminates the
unevaluated merged moment and gives the exact difference
$2\gamma-11/6$. The fixed-rate error estimate in the preceding work
alone would not establish either formula.

\subsection{A quartic example and an evaluated polygamma constant}

For linear rates $a_i=c_i\epsilon$, $c_1<\cdots<c_d$, define the
finite Laurent coefficient
\begin{equation}\label{hc:alpha}
 \alpha_{mi}(\boldsymbol c)=\frac1{c_i-c_m}
 [\epsilon^1]\left\{
   \prod_{k<m}h((c_k-c_i)\epsilon)
   \prod_{\substack{k>m\\k\ne i}}f((c_k-c_i)\epsilon)
 \right\}.
\end{equation}
Then the constant-order part of the complete counterterm is
\begin{equation}\label{hc:linear-constant}
 [\epsilon^0]\mathcal A_{\boldsymbol0}(\epsilon\boldsymbol c)
 =\sum_{m<i}\alpha_{mi}(\boldsymbol c)
               \{\log\epsilon+\log(c_i-c_m)\}.
\end{equation}
In particular, its log-free constant is the finite sum with
$\log(c_i-c_m)$. It requires only \eqref{hc:hseries} through degree
$d-2$. For rational rates every coefficient is a rational polynomial
in $\gamma,\zeta(2),\ldots,\zeta(d-1)$; its formal weight is $d-1$
when $\gamma$ has weight one and $\zeta(k)$ has weight $k$.
Indeed, a coefficient multiplying $z^j$ in either germ has weight
$j+1$, including the pole coefficient at $j=-1$; there are $d-2$
germ factors and their exponents sum to one in \eqref{hc:alpha}.
This is a statement about the produced expressions, not an assertion
of arithmetic independence.

For the quartic rates $(0,1,2,3)$ the six coefficients are
\begin{center}
\begin{tabular}{ccl}
\hline
$(m,i)$ & $c_i-c_m$ & $\alpha_{mi}$\\
\hline
$(1,2)$ & $1$ & $-3\gamma\zeta(2)+\tfrac92\zeta(3)$\\
$(1,3)$ & $2$ & $0$\\
$(1,4)$ & $3$ & $\gamma\zeta(2)-\tfrac32\zeta(3)$\\
$(2,3)$ & $1$ & $\gamma\zeta(2)+4\zeta(3)$\\
$(2,4)$ & $2$ & $2\gamma\zeta(2)-\tfrac92\zeta(3)$\\
$(3,4)$ & $1$ & $5\gamma\zeta(2)$\\
\hline
\end{tabular}
\end{center}
Consequently, with $Q_4=Q_{0,0,0,0}$,
\begin{align}
 \CT_{\epsilon}\mathcal C_{0,0,0,0}(0,\epsilon,2\epsilon,3\epsilon)
 =Q_4
 &+\gamma\zeta(2)(2\log2+\log3)\notag\\
 &-\frac32\zeta(3)(3\log2+\log3).
 \label{hc:quartic}
\end{align}
The coefficient of $\epsilon^0\log\epsilon$ is
$6\gamma\zeta(2)+5\zeta(3)/2$. A rate difference removes $Q_4$:
\begin{align}
 &\CT_{\epsilon}\mathcal C_{0,0,0,0}(0,\epsilon,2\epsilon,4\epsilon)
 -\CT_{\epsilon}\mathcal C_{0,0,0,0}(0,\epsilon,2\epsilon,3\epsilon)
 \notag\\
 &\qquad=\gamma\zeta(2)\left(\frac72\log2+\log3\right)
       -\frac{\zeta(3)}6(\log2+7\log3).
 \label{hc:quartic-difference}
\end{align}

Repeated poles add the rational primitive terms in
\eqref{hc:counter}; retaining only its logarithmic terms would give
an incorrect finite constant. For example,
\begin{equation}\label{hc:middle-derivative}
 \boxed{\quad
 \CT_{\epsilon}\FP_x\int_0^1
   \gamma_0(x)\gamma_0'(\{x+\epsilon\})
                         \gamma_0(\{x+2\epsilon\})\,\dd x
 =2\gamma\zeta(2)-\zeta(3)\log2.
 \quad}
\end{equation}
To check every term, for $\boldsymbol r=(0,1,0)$ and rates $(0,1,2)$,
the three pairs in \eqref{hc:counter} have constant-order contributions
\[
 (m,i)=(1,2):\ -3\zeta(3)\log\epsilon+\zeta(3),\qquad
 (1,3):\ -\zeta(3)(\log\epsilon+\log2),\qquad
 (2,3):\ 4\zeta(3)\log\epsilon.
\]
The nonlogarithmic $\zeta(3)$ in the first pair comes from its double
pole. Hence the finite discrepancy from the merged moment is
$\zeta(3)(1-\log2)$. The merged value follows directly from the
fundamental theorem of calculus before taking the cutoff constant:
\[
 Q_{0,1,0}=\frac13\left\{\gamma^3-
                 \CT_{x\downarrow0}\gamma_0(x)^3\right\}
          =2\gamma\zeta(2)-\zeta(3),
\]
since the constant coefficient of
$(x^{-1}+\gamma-\zeta(2)x+\zeta(3)x^2+\cdots)^3$ is
$\gamma^3-6\gamma\zeta(2)+3\zeta(3)$. Adding the discrepancy proves
\eqref{hc:middle-derivative} without an unevaluated global moment.

\subsection{Verification and scope of the result}

The supplied \src{code/verify\_collisions.py} checks both harmonic
hierarchy laws by exact coefficient extraction for every
$1\le q<p\le12$ and verifies the quartic expressions. Its numerical
calculation independently integrates the original periodic digamma
products, using local Taylor subtraction at each separated seam. The
merged moments are obtained from a local Laurent expansion and ordinary
quadrature away from zero. The collision formula itself is not used
in either quadrature. At 70-digit working precision, selected residuals
$\mathcal C-\mathcal A-Q$ are
\begin{center}
\begin{tabular}{lrr}
\hline
Configuration & $\epsilon=10^{-2}$ & $\epsilon=10^{-3}$\\
\hline
$(0,\epsilon,2\epsilon,3\epsilon)$
 & $-1.468167310\cdot10^{-1}$ & $-1.501685165\cdot10^{-2}$\\
$(0,\epsilon^2,\epsilon)$
 & $-1.288859769\cdot10^{-2}$ & $-1.210723679\cdot10^{-3}$\\
$(0,\epsilon^3,\epsilon^2)$
 & $-1.202913331\cdot10^{-4}$ & $-1.202065558\cdot10^{-6}$\\
$(0,\epsilon,\epsilon+\epsilon^2)$
 & $-1.268176013\cdot10^{-2}$ & $-1.208676388\cdot10^{-3}$\\
\hline
\end{tabular}
\end{center}
These residuals are consistent with a holomorphic remainder, whose
first change is of the order of the largest shift. They are not claimed
to be interval certificates. A separate repeated-pole quadrature in
\src{code/verify\_polygamma\_collisions.py} checks
\eqref{hc:middle-derivative} using second-order local subtraction;
the corresponding residuals are approximately $0.02934754$ and
$0.00224172$. Exact coefficient generation for that example and
additional derivative placements is supplied in
\src{code/derive\_polygamma\_collisions.py}.

The proof of unrestricted multiplicity, argument-derivative order, and
analytic collision rates is the telescoping and contour argument above.
The numerical evidence checks its implementation on specified cases.
The prior fixed-rate triple discrepancy and fixed-collision spectral
completion remain credited to \cite{RepoMixed}. Positive Stieltjes
indices introduce logarithmic singular germs rather than meromorphic
ones; a corresponding exact geometric formula would require an
additional argument and is not asserted here. The unresolved extension
is precisely formulated in the research agenda.
